\documentclass[11pt,letterpaper]{article}

\usepackage[margin=1.5in]{geometry}
\usepackage[T1]{fontenc}
\usepackage[utf8]{inputenc}
\usepackage{lmodern}
\usepackage{microtype}
\usepackage{amsmath,amssymb,amsthm,mathtools,bm}
\usepackage[numbers]{natbib}
\usepackage[hidelinks]{hyperref}
\usepackage{booktabs}
\usepackage{enumitem}

\hypersetup{
  pdftitle={Contraction Metrics and Tangent-Space Optimal Control},
  pdfauthor={AffineDrift},
  pdfsubject={Nonlinear control, contraction theory, and local optimal control},
  pdfkeywords={contraction metrics, tangent dynamics, LQR, DDP, Riccati, robotics, biomechanics}
}

\newtheorem{theorem}{Theorem}[section]
\newtheorem{proposition}[theorem]{Proposition}
\newtheorem{definition}[theorem]{Definition}
\newtheorem{remark}[theorem]{Remark}

\newcommand{\R}{\mathbb{R}}
\newcommand{\dd}{\mathrm{d}}
\newcommand{\norm}[1]{\left\lVert #1 \right\rVert}
\newcommand{\M}{\mathbf{M}}
\newcommand{\Smat}{\mathbf{S}}
\newcommand{\A}{\mathbf{A}}
\newcommand{\B}{\mathbf{B}}
\newcommand{\Q}{\mathbf{Q}}
\newcommand{\W}{\mathbf{W}}
\newcommand{\x}{\mathbf{x}}
\newcommand{\uvec}{\mathbf{u}}
\newcommand{\dx}{\delta\mathbf{x}}
\newcommand{\du}{\delta\mathbf{u}}

\title{Contraction Metrics and Tangent-Space Optimal Control\\\large A Repaired Working Draft for Ongoing Development}
\author{AffineDrift Working Notes}
\date{\today}

\begin{document}
\maketitle

\begin{abstract}
This document develops a unified viewpoint for nonlinear stability and local optimal control through tangent-space dynamics. The central idea is that contraction analysis and Riccati-based trajectory optimization act on the same infinitesimal perturbations. This repaired draft is intentionally organized as a writing-ready technical manuscript: each section includes (i) an introductory qualitative framing, (ii) an ``In Laymans Terms'' explanation, and (iii) a ``Geometric Intuition'' explanation, followed by detailed equations and interpretation. The article is designed to be edited further into textbook or paper form.
\end{abstract}

\tableofcontents
\newpage

\section{Roadmap and Scope}

\subsection{Introduction and Qualitative Purpose}
The purpose of this manuscript is to make one technical claim precise: local optimal control and contraction-based stability certificates can be interpreted using the same tangent-space objects. Rather than presenting independent threads, we build a common language first (variational dynamics and metrics), then show how Lyapunov-style and Riccati-style arguments become different projections of the same geometry.

The practical purpose of the calculations is equally important. In high-dimensional systems, equations are not written only for proof elegance; they are written to answer engineering questions: Which perturbations grow fastest? Which control channels are most effective? Which local model remains trustworthy over a finite horizon? Each derivation below is accompanied by a qualitative statement of what that equation buys you in design or diagnostics.

\subsection{In Laymans Terms}
This article is about how to keep complex motions both \emph{good} and \emph{repeatable}. ``Good'' means performance with respect to a goal (for example, tracking, energy use, impact quality). ``Repeatable'' means small disturbances do not derail the motion. The main message is that both goals can be addressed by studying how tiny errors evolve over time.

\subsection{Geometric Intuition}
A trajectory is a curve on a state-space surface. Around each point on that curve, there is a tangent plane that captures first-order motion of nearby trajectories. All local stability and local optimality calculations are happening in these moving tangent planes. A metric is a local ruler on those planes; changing the metric changes which perturbation directions are considered large or small.

\subsection{Article Structure}
The manuscript is organized into ten technical sections and one conclusion:
\begin{enumerate}[label=\arabic*.]
\item Contraction theory fundamentals.
\item Tangent-space dynamics and local linearization.
\item LQR and Differential Dynamic Programming (DDP).
\item Stability-optimality duality through Riccati objects.
\item Continuous-time and infinite-horizon limits.
\item Coordinate-free and Riemannian interpretation.
\item Contraction-constrained computational design.
\item Application patterns in biomechanics and robotics.
\item Implementation checklist for practical use.
\item Consolidated conclusion and forward path.
\end{enumerate}

\section{Contraction Theory Primer}

\subsection{Introduction and Qualitative Purpose}
Contraction theory answers a trajectory-to-trajectory question: if two trajectories start near each other, do they converge exponentially? This is a stronger statement than local equilibrium stability when the behavior of interest is a moving trajectory rather than a fixed point.

Qualitatively, contraction calculations provide a disturbance-forgetting certificate. If a region is contracting, the system attenuates uncertainty and initialization mismatch at a quantifiable rate. In practice this gives a robust foundation for planning and control over families of trajectories rather than single nominal motions.

\subsection{In Laymans Terms}
Imagine running the same robot motion twice with tiny initial differences. If the system is contracting, those two runs quickly become almost identical. That means small mistakes do not snowball.

\subsection{Geometric Intuition}
Contraction checks whether tangent vectors between nearby trajectories shrink under the flow. If all tangent vectors shrink in a chosen metric, geodesic distance between trajectories also shrinks. So contraction is local shrinking that integrates into global convergence behavior on a region.

\subsection{Differential Dynamics and Metric Condition}
Consider
\begin{equation}
\dot{\x} = f(\x,t), \quad \x \in \R^n.
\label{eq:nonlinear-system}
\end{equation}
The variational dynamics are
\begin{equation}
\dot{\dx} = A(\x,t)\,\dx, \quad A(\x,t) := \frac{\partial f}{\partial x}(\x,t).
\label{eq:variational-dynamics}
\end{equation}
Let $\M(\x,t) \succ 0$ be a smooth metric and define $V = \dx^\top \M\dx$. Then
\begin{equation}
\dot{\M}(\x,t)
:= \frac{\partial \M}{\partial t}(\x,t)
 + \sum_{i=1}^{n}\frac{\partial \M}{\partial x_i}(\x,t)\,f_i(\x,t),
\label{eq:metric-total-derivative}
\end{equation}
where $\dot{\M}$ is the total derivative of the metric along the flow of \eqref{eq:nonlinear-system}. Using this definition,
\begin{equation}
\dot V = \dx^\top\!\left(A^\top\M + \M A + \dot{\M}\right)\!\dx.
\label{eq:metric-derivative}
\end{equation}
A sufficient contraction condition with rate $\lambda>0$ is
\begin{equation}
A^\top\M + \M A + \dot{\M} \preceq -2\lambda\M.
\label{eq:contraction-lmi}
\end{equation}

\begin{theorem}[Standard contraction implication]
Assume $\mathcal{D}$ is forward invariant and geodesically convex in the metric $\M(\x,t)$. Assume also that there exist constants $0<\underline{m}\leq \overline{m}<\infty$ such that
\begin{equation}
\underline{m}I \preceq \M(\x,t) \preceq \overline{m}I
\quad \text{for all } (\x,t)\in\mathcal{D}\times[0,\infty).
\label{eq:metric-bounds}
\end{equation}
If \eqref{eq:contraction-lmi} holds uniformly on $\mathcal{D}$, then for any two trajectories $\x_1(t),\x_2(t)$ in $\mathcal{D}$,
\begin{equation}
d_{\M,t}\!\left(\x_1(t),\x_2(t)\right)
\le e^{-\lambda t}\,d_{\M,0}\!\left(\x_1(0),\x_2(0)\right).
\label{eq:geodesic-decay}
\end{equation}
Moreover, Euclidean separation satisfies
\begin{equation}
\norm{\x_1(t)-\x_2(t)}
\le \sqrt{\frac{\overline{m}}{\underline{m}}}\,e^{-\lambda t}
\norm{\x_1(0)-\x_2(0)}.
\label{eq:euclidean-decay}
\end{equation}
\end{theorem}

\subsection{Why This Matters for Design}
The condition in \eqref{eq:contraction-lmi} provides an engineering lever: select or optimize a metric $\M$ that makes the symmetric part of the Jacobian sufficiently negative. This creates a direct bridge to convex and non-convex metric synthesis methods discussed later.

\section{Tangent-Space Dynamics and Local Linearization}

\subsection{Introduction and Qualitative Purpose}
Most nonlinear control algorithms repeatedly linearize along a nominal trajectory. This is not merely approximation convenience; it is the correct infinitesimal description of how perturbations propagate. The qualitative role of tangent-space equations is to separate \emph{what the nominal motion does} from \emph{how deviations evolve around it}.

\subsection{In Laymans Terms}
At each instant, pretend the nonlinear system is locally linear around what it is currently doing. Use that local linear model to predict whether tiny errors will grow or shrink.

\subsection{Geometric Intuition}
A tangent vector is a local displacement direction. Linearization gives a map from one tangent space to the next along the trajectory. In discrete time this map is a matrix $A_t$; in continuous time it is a Jacobian field $A(x,t)$.

\subsection{Control-Affine Model and Discrete Perturbation System}
For control-affine dynamics,
\begin{equation}
\dot x = f(x,t) + B(x,t)u,
\label{eq:affine-dynamics}
\end{equation}
the continuous-time variational model is
\begin{equation}
\dot{\dx} = A(t)\dx + B(t)\du.
\label{eq:ct-variational-control-affine}
\end{equation}
For sampled-data implementations with sampling period $h$ and zero-order-hold input, define the discrete flow map
\begin{equation}
x_{t+1} = f_d(x_t,u_t) := \Phi_h(x_t,u_t),
\label{eq:discrete-flow-map}
\end{equation}
where $\Phi_h$ is the time-$h$ solution operator induced by \eqref{eq:affine-dynamics}. Linearization along $(x_t^\star,u_t^\star)$ then yields
\begin{equation}
\dx_{t+1} = A_t\,\dx_t + B_t\,\du_t,
\label{eq:discrete-perturbation}
\end{equation}
with Jacobians
\begin{equation}
A_t = \left.\frac{\partial f_d}{\partial x}\right|_{(x_t^\star,u_t^\star)},
\qquad
B_t = \left.\frac{\partial f_d}{\partial u}\right|_{(x_t^\star,u_t^\star)}.
\label{eq:ab-jacobians}
\end{equation}

\subsection{Qualitative Interpretation of the Matrices}
$A_t$ is local state sensitivity; it tells which perturbation modes naturally amplify or contract. $B_t$ is local control authority; it tells which perturbation directions can be corrected efficiently. Many practical failures in optimization-based control are actually poor conditioning of these local maps, not failure of the nominal planner itself.

\section{LQR and Differential Dynamic Programming}

\subsection{Introduction and Qualitative Purpose}
LQR and DDP solve local optimal correction problems in tangent coordinates. Their qualitative role is to convert linearized sensitivity information into a feedback law that trades error reduction against control effort.

In the trajectory setting, DDP repeatedly refreshes the local model and local quadratic cost. That loop can be viewed as ``fit local geometry, solve local optimization, update base curve.''

\subsection{In Laymans Terms}
Take the current best motion guess. Zoom in so the system looks linear and the cost looks quadratic. Solve that simpler problem exactly. Update the guess and repeat.

\subsection{Geometric Intuition}
The Riccati matrix is a time-varying quadratic form on perturbations. Geometrically, each step assigns local curvature to the value landscape in tangent space. Steep curvature means errors in that direction are expensive and should be corrected aggressively.

\subsection{Finite-Horizon LQR Subproblem}
Given \eqref{eq:discrete-perturbation}, minimize
\begin{equation}
J = \frac{1}{2}\dx_T^\top Q_T\dx_T
 + \frac{1}{2}\sum_{t=0}^{T-1}
\left(\dx_t^\top Q_t\dx_t + \du_t^\top R_t\du_t\right),
\label{eq:lqr-cost}
\end{equation}
with $Q_t\succeq0$, $R_t\succ0$. The optimal local law is
\begin{equation}
\du_t = -K_t\dx_t,
\label{eq:feedback-law}
\end{equation}
where
\begin{equation}
\begin{aligned}
S_T &= Q_T,\\
K_t &= (R_t + B_t^\top S_{t+1}B_t)^{-1} B_t^\top S_{t+1}A_t,\\
S_t &= Q_t + A_t^\top S_{t+1}A_t \\
&\quad - A_t^\top S_{t+1}B_t(R_t + B_t^\top S_{t+1}B_t)^{-1}B_t^\top S_{t+1}A_t.
\end{aligned}
\label{eq:discrete-riccati}
\end{equation}
An equivalent identity, useful for Lyapunov and contraction arguments, is
\begin{equation}
S_t = Q_t + K_t^\top R_tK_t + (A_t-B_tK_t)^\top S_{t+1}(A_t-B_tK_t).
\label{eq:riccati-equivalent}
\end{equation}

\subsection{DDP Interpretation}
DDP augments this with second-order model updates and forward rollouts, but the local stabilizing core remains the same: each backward pass computes a curvature object $S_t$ and a gain $K_t$ that shape local error dynamics.

\section{Stability-Optimality Duality}

\subsection{Introduction and Qualitative Purpose}
This section states the central bridge: under linear-quadratic assumptions, the Riccati matrix from optimal control can serve as a contraction metric for closed-loop tangent dynamics. The qualitative benefit is unification: one matrix sequence can certify both performance (value) and local convergence (stability).

\subsection{In Laymans Terms}
The same matrix used by LQR to score future cost can also be used as a ruler proving that small errors shrink under the feedback law.

\subsection{Geometric Intuition}
Optimal control builds a local quadratic ``cost bowl''; contraction checks shrinking under a metric bowl. Duality appears when those bowls coincide for closed-loop perturbation dynamics.

\subsection{Main Statement (Discrete-Time Form)}
Closed-loop perturbations satisfy
\begin{equation}
\dx_{t+1} = (A_t-B_tK_t)\dx_t.
\label{eq:closed-loop-perturbation}
\end{equation}
Define $V_t = \dx_t^\top S_t\dx_t$ with $S_t$ from \eqref{eq:discrete-riccati}. Then
\begin{equation}
V_{t+1} - V_t
= -\dx_t^\top(Q_t + K_t^\top R_t K_t)\dx_t.
\label{eq:value-drop}
\end{equation}
Hence, if there exist constants $\alpha,\underline{s},\overline{s}>0$ such that for all $t$,
\begin{equation}
Q_t + K_t^\top R_t K_t \succeq \alpha I,
\qquad
\underline{s}I \preceq S_t \preceq \overline{s}I,
\label{eq:duality-bounds}
\end{equation}
and $0<\alpha<\overline{s}$, one obtains exponential decay
\begin{equation}
V_{t+1} \le \left(1-\frac{\alpha}{\overline{s}}\right)V_t.
\label{eq:discrete-exp-decay}
\end{equation}
and therefore
\begin{equation}
\norm{\dx_t}^2
\le
\frac{\overline{s}}{\underline{s}}
\left(1-\frac{\alpha}{\overline{s}}\right)^t
\norm{\dx_0}^2.
\label{eq:state-norm-decay}
\end{equation}
This is a contraction-style inequality in the metric induced by $S_t$, with an explicit Euclidean-norm consequence.

\subsection{Design Implication}
Instead of treating trajectory optimization and contraction verification as separate pipelines, one can interpret Riccati outputs as candidate metrics, then verify robustness margins under model error and constraints.

\section{Continuous-Time and Infinite-Horizon Limits}

\subsection{Introduction and Qualitative Purpose}
Continuous-time analysis clarifies rate interpretations and provides links to classical control design equations. Infinite-horizon limits further reduce transient complexity and expose steady-state structure used in synthesis and certification.

\subsection{In Laymans Terms}
If the task and dynamics do not change too much over time, the ``best local ruler'' and ``best local feedback law'' settle to steady matrices.

\subsection{Geometric Intuition}
The differential Riccati equation evolves a time-varying metric tensor on perturbations. In the infinite-horizon limit that tensor can converge to a fixed geometric object, giving constant local curvature of value and contraction.

\subsection{Differential and Algebraic Riccati Equations}
For
\begin{equation}
\dot{\dx} = A(t)\dx + B(t)\du,
\label{eq:ct-perturbation}
\end{equation}
with local quadratic density $\frac12\dx^\top Q\dx + \frac12\du^\top R\du$, the backward differential Riccati equation is
\begin{equation}
-\dot S = Q + A^\top S + SA - SBR^{-1}B^\top S.
\label{eq:dre}
\end{equation}
For time-invariant data, the infinite-horizon matrix $S\succeq0$ solves the algebraic Riccati equation (ARE)
\begin{equation}
Q + A^\top S + SA - SBR^{-1}B^\top S = 0.
\label{eq:are}
\end{equation}
Closed-loop dynamics $\dot{\dx} = (A-BK)\dx$ with $K=R^{-1}B^\top S$ satisfy
\begin{equation}
\dot V = -\dx^\top(Q+K^\top RK)\dx,
\quad V=\dx^\top S\dx,
\label{eq:ct-value-drop}
\end{equation}
which is the continuous-time analogue of the duality argument.

\section{Coordinate-Free and Riemannian Formulation}

\subsection{Introduction and Qualitative Purpose}
Many systems of interest naturally live on manifolds (orientation, rigid body pose, constrained linkage coordinates). A coordinate-free perspective keeps statements invariant under reparameterization and avoids misleading artifacts from poor coordinates.

Qualitatively, this section explains how to interpret local quadratic forms as metric tensors, and why pullbacks are the correct way to move task penalties into configuration coordinates.

\subsection{In Laymans Terms}
Some state spaces are curved, not flat. The math still works if we replace ordinary dot products with local curved-space rulers.

\subsection{Geometric Intuition}
A Riemannian metric assigns an inner product to each tangent space. Control and estimation penalties are then geometric lengths/energies of perturbation vectors. Coordinate changes modify matrices but not the underlying geometric object.

\subsection{Metric Tensor and Pullback}
Let $(\mathcal{X},g)$ be a Riemannian state manifold. The local perturbation energy is
\begin{equation}
\norm{\dx}_{g}^2 = g_x(\dx,\dx).
\label{eq:riemannian-energy}
\end{equation}
In coordinates, $g_x$ is represented by $M(x)\succ0$ and \eqref{eq:riemannian-energy} becomes $\dx^\top M(x)\dx$.

For a task map $y=h(x)$ with Jacobian $J(x)$ and task metric $W(y)\succ0$, the pullback metric is
\begin{equation}
M(x) = J(x)^\top W(h(x))J(x).
\label{eq:pullback-metric}
\end{equation}
If $J(x)$ has full column rank, \eqref{eq:pullback-metric} is positive definite and defines a true Riemannian metric. If $J(x)$ is rank-deficient, \eqref{eq:pullback-metric} is positive semidefinite; in that case a common regularization is
\begin{equation}
M_\varepsilon(x) = J(x)^\top W(h(x))J(x) + \varepsilon I,
\quad \varepsilon>0,
\label{eq:pullback-regularized}
\end{equation}
or an equivalent nullspace completion. These constructions preserve the same physical interpretation: penalize error where the task lives, then pull that penalty back to generalized coordinates.

\subsection{Coordinate Change Consistency}
Under a smooth reparameterization $z=\phi(x)$ with Jacobian $T=\partial x/\partial z$, metric matrices transform as
\begin{equation}
M_z(z) = T(z)^\top M_x(x(z))T(z),
\label{eq:metric-transform}
\end{equation}
ensuring geometric statements remain invariant.

\section{Contraction-Constrained Design and Computation}

\subsection{Introduction and Qualitative Purpose}
This section focuses on synthesis: how to compute metrics and feedback laws that are not only performance-optimal locally but also satisfy explicit contraction margins. The qualitative theme is to turn robustness requirements into optimization constraints.

\subsection{In Laymans Terms}
Do not only ask for a controller that performs well in simulation. Also require a mathematical shrinking margin for small errors, and solve for both at once.

\subsection{Geometric Intuition}
Optimization chooses a local ruler and local control directions so the flow bends perturbation vectors back toward the nominal trajectory while still achieving task objectives.

\subsection{Representative Optimization Programs}
\paragraph{LTI case (convex).}
Given $(A,B)$, one can search for $M\succ0$ and gain $K$ such that
\begin{equation}
(A-BK)^\top M + M(A-BK) \preceq -2\lambda M.
\label{eq:lti-contraction-synthesis}
\end{equation}
With suitable variable substitutions, this can be posed as an LMI and solved using semidefinite programming \citep{Boyd1994}.

\paragraph{Trajectory case (time-varying).}
Given nominal $(A_t,B_t)$ from rollouts, solve stage-wise or horizon-coupled programs for $(S_t,K_t)$ with additional inequalities enforcing a minimum contraction rate.

\paragraph{Nonlinear case (sampled).}
Sample points along candidate trajectories; enforce discretized versions of \eqref{eq:contraction-lmi}; optimize a parameterized metric ansatz (for example neural or polynomial) using autodiff and constrained optimization.

\subsection{Practical Numerical Workflow}
\begin{enumerate}[label=\arabic*.]
\item Generate a nominal trajectory and linearization data $(A_t,B_t)$.
\item Solve a baseline Riccati recursion for $(K_t,S_t)$.
\item Evaluate contraction residuals from \eqref{eq:contraction-lmi} or \eqref{eq:lti-contraction-synthesis}.
\item If margins are weak, re-optimize weights/metrics and repeat.
\item Validate decay rates by Monte Carlo perturbation rollouts.
\end{enumerate}

\section{Applications in Biomechanics and Robotics}

\subsection{Introduction and Qualitative Purpose}
Applications motivate why this unification matters. In biomechanics and robotics, motions are high-dimensional, contact-sensitive, and strongly coupled. A useful theory must say both how to optimize movement and why nearby executions remain coherent under disturbances.

\subsection{In Laymans Terms}
For a golf swing, manipulation task, or gait, you want two properties: high performance and repeatability. The tangent-space framework gives a way to tune both with the same local models.

\subsection{Geometric Intuition}
Complex multi-joint motion is a trajectory on a high-dimensional manifold. Muscle groups, actuators, and contact constraints shape which tangent directions are controllable. Metrics encode which deviation directions matter most for task success.

\subsection{Biomechanics Pattern}
For multibody dynamics
\begin{equation}
H(q)\ddot q + C(q,\dot q)\dot q + g(q) = \tau + J_c(q)^\top\lambda,
\label{eq:multibody}
\end{equation}
where $H(q)$ is the inertia matrix (using $H$ avoids notational conflict with contraction metrics $M$), the perturbation model reveals sensitivity to inertial coupling and contact terms. Metric shaping can prioritize perturbation suppression in joint combinations associated with outcome variables (club-head speed, impact orientation, etc.).

\subsection{Robotics Pattern}
In operational space control, a task Jacobian $J(q)$ maps joint perturbations to task perturbations. Pullback metrics from task coordinates (Equation \eqref{eq:pullback-metric}) provide a direct way to encode anisotropic task priorities while preserving coordinate consistency.

\subsection{What to Measure in Experiments}
\begin{itemize}
\item Local decay rate of perturbation norm in chosen metric.
\item Cost-to-go decrease consistency with Riccati predictions.
\item Sensitivity of contraction margin under parameter mismatch.
\item Basin size estimates from perturbation ensembles.
\end{itemize}

\section{Implementation Checklist for Ongoing Work}

\subsection{Introduction and Qualitative Purpose}
This section is a practical editing and validation scaffold so the document can evolve without losing mathematical consistency. The qualitative goal is repeatable development: every new derivation should map to a tangible diagnostic and a clear geometric interpretation.

\subsection{In Laymans Terms}
Use this as a quality checklist while extending the paper. Every new formula should answer ``what does it mean physically?'' and ``how would we test it?''

\subsection{Geometric Intuition}
Treat every section as a map from equations to geometry: identify the tangent object, identify the metric, identify the contraction or optimization claim, and identify the coordinate-invariance statement.

\subsection{Authoring Checklist}
\begin{enumerate}[label=\arabic*.]
\item Keep one nominal trajectory notation throughout a section.
\item State dimensions of all new matrices at first use.
\item Pair each theorem with an engineering interpretation paragraph.
\item Add one ``failure mode'' paragraph where assumptions can break.
\item Include a reproducibility note for any numerical claim.
\item Preserve the three explanatory blocks in every section.
\end{enumerate}

\subsection{Computation Checklist}
\begin{enumerate}[label=\arabic*.]
\item Verify linearization accuracy against finite differences.
\item Track condition numbers of $R_t + B_t^\top S_{t+1}B_t$.
\item Log contraction residual eigenvalues along trajectories.
\item Validate predicted decay rates in closed-loop rollouts.
\item Stress-test against model perturbations and contact changes.
\end{enumerate}

\section{Conclusion and Forward Path}

\subsection{Introduction and Qualitative Purpose}
The repaired draft establishes a coherent narrative where contraction and local optimal control share a tangent-space foundation. The immediate purpose is to provide a stable base for deeper theorem proofs, richer case studies, and code-linked reproducibility appendices.

\subsection{In Laymans Terms}
The key message is simple: the same local mathematics can help you make motions better and more reliable. This document is now clean enough to keep building on.

\subsection{Geometric Intuition}
A trajectory, its tangent perturbations, and a metric-defined ruler form a minimal geometric toolkit. Stability and optimality are two questions asked of that same toolkit.

\subsection{Near-Term Expansion Targets}
\begin{enumerate}[label=\arabic*.]
\item Add a full proof section for discrete and continuous duality conditions.
\item Add one complete worked example with numeric matrices and plots.
\item Add a short appendix on implementation details (autodiff, SDP, and sampling).
\item Add explicit robustness margins under bounded modeling error.
\end{enumerate}

\section*{References}
\addcontentsline{toc}{section}{References}
\bibliographystyle{plainnat}
\begin{thebibliography}{99}

\bibitem[Lohmiller and Slotine(1998)]{LohmillerSlotine1998}
Lohmiller, W. and Slotine, J. J. E. (1998).
On contraction analysis for non-linear systems.
\emph{Automatica}, 34(6), 683--696.

\bibitem[Jouffroy and Slotine(2004)]{JouffroySlotine2004}
Jouffroy, J. and Slotine, J. J. E. (2004).
Methodological remarks on contraction theory.
In \emph{Proceedings of the 43rd IEEE Conference on Decision and Control}, 2537--2543.

\bibitem[Manchester and Slotine(2017)]{ManchesterSlotine2017}
Manchester, I. R. and Slotine, J. J. E. (2017).
Control contraction metrics and model-based control of nonlinear systems.
\emph{IEEE Transactions on Automatic Control}, 62(9), 4506--4521.

\bibitem[Bullo(2024)]{BulloContraction}
Bullo, F. (2024).
\emph{Contraction Theory for Dynamical Systems}.
Manuscript.

\bibitem[Khalil(2002)]{Khalil2002}
Khalil, H. K. (2002).
\emph{Nonlinear Systems} (3rd ed.).
Prentice Hall.

\bibitem[Sastry(1999)]{Sastry1999}
Sastry, S. (1999).
\emph{Nonlinear Systems: Analysis, Stability, and Control}.
Springer.

\bibitem[Slotine and Li(1991)]{SlotineLi1991}
Slotine, J. J. E. and Li, W. (1991).
\emph{Applied Nonlinear Control}.
Prentice Hall.

\bibitem[Isidori(1995)]{Isidori1995}
Isidori, A. (1995).
\emph{Nonlinear Control Systems} (3rd ed.).
Springer.

\bibitem[Mayne(1966)]{Mayne1966}
Mayne, D. Q. (1966).
A second-order gradient method for determining optimal trajectories of non-linear discrete-time systems.
\emph{International Journal of Control}, 3(1), 85--95.

\bibitem[Todorov and Jordan(2002)]{TodorovJordan2002}
Todorov, E. and Jordan, M. I. (2002).
Optimal feedback control as a theory of motor coordination.
\emph{Nature Neuroscience}, 5(11), 1226--1235.

\bibitem[Tedrake(2009)]{Tedrake2009}
Tedrake, R. (2009).
LQR-Trees: Feedback motion planning on sparse randomized trees.
In \emph{Robotics: Science and Systems}.

\bibitem[Spong et~al.(2006)]{Spong2006}
Spong, M. W., Hutchinson, S., and Vidyasagar, M. (2006).
\emph{Robot Modeling and Control}.
Wiley.

\bibitem[Khatib(1987)]{Khatib1987}
Khatib, O. (1987).
A unified approach for motion and force control of robot manipulators: The operational space formulation.
\emph{IEEE Journal on Robotics and Automation}, 3(1), 43--53.

\bibitem[Boyd et~al.(1994)]{Boyd1994}
Boyd, S., El Ghaoui, L., Feron, E., and Balakrishnan, V. (1994).
\emph{Linear Matrix Inequalities in System and Control Theory}.
SIAM.

\end{thebibliography}

\end{document}
